%====================================================================
% Trace‑Driven Multi‑Service Workflow Testing for Microservice Architectures
%====================================================================
\documentclass[conference]{IEEEtran}

% ---------- Packages ------------
\usepackage{graphicx}
\usepackage{listings}
\usepackage{hyperref}
\usepackage{amsmath}
\usepackage{enumitem}
\usepackage{xcolor}
\usepackage{booktabs}
\usepackage{tabularx}
\usepackage{array}
\usepackage{algorithm}
\usepackage{algorithmicx}
\usepackage{algpseudocode}
\usepackage{url}
\usepackage{doi} 
\setlist[itemize]{leftmargin=*,topsep=2pt,itemsep=2pt}
\newenvironment{compactdesc}{\begin{description}[leftmargin=0pt,style=nextline]}{\end{description}}

% ---------- Metadata ------------
\title{AI‑Enhanced Trace‑Driven API Testing with Comprehensive Microservice Validation}

\author{\IEEEauthorblockN{Tingshuo Miao}\IEEEauthorblockA{Department of Computer Science\\Baylor University, USA\\Email: tingshuo\_miao2@baylor.edu}}

%====================================================================
\begin{document}
\maketitle

% ---------------- Abstract ----------------
\begin{abstract}
Microservice architectures present unique testing challenges where failures often emerge through complex service interactions rather than individual API faults. Traditional workflow-based testing approaches face fundamental execution challenges: replaying multi-step sequences as separate client requests fails to capture ephemeral in-memory data dependencies and creates misleading test oracles. This paper presents an AI-enhanced, trace-driven extension to RESTest that employs \emph{Root API Mode} testing: (i) extracting publicly accessible API endpoints from distributed traces, (ii) generating diverse test cases targeting these entry points using smart input fetching and large language model synthesis, (iii) performing comprehensive post-execution trace analysis to validate complete microservice interaction chains triggered by single API calls, and (iv) maintaining a persistent Root API Registry that systematically catalogues discovered APIs with their interaction patterns for long-term testing strategy optimization. Our approach integrates \textit{TraceErrorAnalyzer} with LLM-powered failure diagnosis to provide actionable root cause analysis for complex distributed failures. The system combines \textit{SmartInputFetcher} for realistic parameter generation with \textit{AiDrivenLLMGenerator} for exploratory testing, executing tests through enhanced Allure reporting that visualizes complete service interaction hierarchies. This Root API Mode methodology achieves comprehensive microservice validation with superior efficiency and accuracy compared to traditional workflow replication while providing unprecedented failure diagnosis capabilities through AI-enhanced trace analysis.
\end{abstract}

\begin{IEEEkeywords}
Microservices, REST API testing, Root API Mode, distributed tracing, automated test generation, trace analysis, AI-enhanced testing, RESTest.
\end{IEEEkeywords}

%====================================================================
\section{Introduction}
The microservice architectural style has become the \emph{de-facto}
standard for cloud-native software because it promises independent
deployment, elastic scaling, and organisational agility.  Yet this
fine-grained decomposition comes at a price: many faults emerge only
when several services collaborate along an end-to-end workflow, whereas
most automated tools still exercise each API in isolation.

\subsection*{Why is Testing Microservices Hard?}
Industry surveys and mapping studies consistently identify testing as
one of the costliest and most error-prone activities in
microservice-based development
\cite{waseem2020testing,miao2025sms}.  
Three challenges appear throughout the literature:

\begin{enumerate}[leftmargin=*,label=\roman*)]
  \item \textbf{API correctness} — Every service publishes an
        independent REST/gRPC/GraphQL contract whose input domain must be
        explored thoroughly; malformed requests or untested edge cases
        can propagate faults to upstream clients.
  \item \textbf{Workflow integrity} — Real business transactions
        traverse \emph{multiple} services.  
        Integration-relevant failures such as data-flow mismatches,
        missing authentication tokens, or out-of-order requests
        represent a substantial share of production incidents
        \cite{zhou2018microservice}.
  \item \textbf{Resilience under partial failure} — Network partitions,
        container restarts, and eventual-consistency delays mandate
        fault-injection (a.k.a.\ chaos) testing, yet fewer than one
        third of automated tools provide built-in resilience validation
        \cite{waseem2020testing}.
\end{enumerate}

A recent \emph{Systematic Mapping Study of Test Generation for
Microservices} \cite{miao2025sms} groups existing work into
\textit{consumer-driven contract}, \textit{search-based input
generation}, and \textit{resilience testing}.  Crucially, it concludes
that no current tool unifies these approaches into a single,
workflow-aware framework—leaving integration faults largely
undetected.

\subsection*{Problem Statement}
\textsc{RESTest}~\cite{martin2019restest} epitomises the state of the
art in automated REST API testing: it generates HTTP requests that
achieve high structural and semantic coverage for \emph{single}
endpoints.  However, \textsc{RESTest} and comparable tools cannot
validate the complex microservice interactions that occur when external clients invoke publicly accessible API endpoints.

\subsection*{Objective}
We aim to extend \textsc{RESTest} to automatically identify and test \emph{root API endpoints}—the externally accessible entry points that trigger comprehensive microservice interactions—while providing deep visibility into the complete service interaction chains through post-execution trace analysis.

\subsection*{Approach Overview}
Our key insight is to exploit \emph{distributed traces} to identify the publicly accessible API endpoints that serve as entry points to complex microservice ecosystems. By analyzing these traces, we extract root APIs and their parameter patterns, generate diverse test variants targeting these entry points, and then perform comprehensive post-execution trace analysis to validate the complete microservice interaction chains. This \emph{Root API Mode} approach avoids the fundamental challenges of workflow replication while providing superior validation coverage through intelligent trace-based analysis. Figure~\ref{fig:dataflow} sketches the resulting pipeline; the methodological details are developed in Section~\ref{sec:method}.


% \textbf{Contributions.}


%====================================================================
\section{Related Work}
\label{sec:related}

\subsection{API-level test generation for Web/microservice APIs}
Model-based and search-based test generators remain the backbone of automated API testing. 
\textsc{RESTest}~\cite{martin2019restest} extracts a test model from OpenAPI and generates black-box tests that maximise structural and HTTP-state coverage for \emph{single} endpoints, with pluggable data generators and Allure reporting. 
EvoMaster~\cite{arcuri2019restful} pushes search-based generation further (white-/black-box) and has since been extended to RPC and GraphQL settings with industrial evaluations.
Complementary fuzzers such as RESTler~\cite{atlidakis2019restler} analyse API specs to infer producer--consumer dependencies and exercise \emph{stateful} request sequences; Pythia~\cite{atlidakis2020pythia} augments grammar-based fuzzing with coverage-guided feedback; foREST~\cite{lin2022forest} introduces tree-based dependency abstractions; and MINER~\cite{lyu2023miner} leverages neural predictors for parameter values. 
RestTestGen~\cite{viglianisi2020resttestgen} offers automated black-box testing from OpenAPI, while recent work targets specific vulnerability classes such as excessive data exposure~\cite{pan2023edefuzz}.

Unlike these tools, which predominantly treat endpoints in isolation or rely on end-to-end sequence replay, our \emph{Root API Mode} anchors generation at externally reachable entry points and validates the \emph{entire} induced interaction chain \emph{post-execution} via traces, thus avoiding oracle pitfalls caused by artificial multi-step replays.

\subsection{Contract testing and integration verification}
Consumer-Driven Contract (CDC) testing mitigates interface drift in independently deployed services. 

Empirical accounts show CDC’s practicality in industry (e.g., a case study by Lehv{\"a} el.~\cite{lehva2019cdc}), and mapping studies highlight its role in microservice ecosystems~\cite{waseem2020testing}. 

CDC complements our approach by stabilising inter-service syntactic and semantic expectations; however, CDC alone does not validate \emph{runtime} cross-service causality or internal orchestration behaviours that we verify through trace-based post-execution analysis.

\subsection{Grey-box and end-to-end testing for microservices}
Beyond black-box API testing, grey-box approaches instrument the runtime to recover inter-service interactions and internal coverage. 
Giamattei et al.’s MACROHIVE~\cite{giamattei2022macrohive} automatically generates and executes tests with monitoring proxies to reveal interaction-level coverage and causal failure chains. 

At the system level, recent benchmarks and artefact surveys facilitate evaluation: Smith el.~\cite{smith2023e2e-benchmark} provides end-to-end testing benchmarks, while Fischer el.~\cite{fischer2024overview} curates widely used microservice systems and datasets for testing/monitoring research. 

Our Root API Mode differs by \emph{eschewing} explicit replay of multi-step workflows: it executes only public entry points and relies on \emph{distributed traces} to validate full chains, thereby preserving ephemeral, in-memory dependencies that grey-box or E2E replays may disrupt.

\subsection{Trace-based diagnosis, observability, and failure analysis}
Distributed tracing has long served causal analysis at scale (e.g., Dapper~\cite{sigelman2010dapper}, Pivot Tracing~\cite{mace2015pivot}). 
For microservices, empirical studies underscore tracing’s centrality for fault analysis and debugging~\cite{zhou2018microservice}. 
A growing body of work localises root causes directly from traces: TraceRCA ~\cite{li2021tracerca} ranks suspect services via anomalous-trace statistics; MicroHECL ~\cite{liu2021microhecl} correlates multi-type anomalies on dynamic call graphs; ImpactTracer~\cite{xie2023impacttracer} builds impact graphs and performs backward search; and recent work explores multi-dimensional trace signals for localisation~\cite{zhang2024tracecontrast}. 
Our \textsc{TraceErrorAnalyzer} operates in this lineage but is tightly integrated into the testing loop: it attaches trace retrieval to each Root API invocation, performs hierarchical failure analysis over the induced call tree, and (optionally) leverages LLMs for actionable diagnosis---all while avoiding the artefacts of workflow replays.

\subsection{Coverage metrics for microservice systems}
Classical API coverage criteria (parameter, operation, status-code, response-body)~\cite{martin2019restest,arcuri2019restful} do not capture the breadth of \emph{induced} interactions across services. 
Recent work proposes endpoint-centric E2E and API coverage metrics, with tooling that maps tests to endpoints across microservices and visualises uncovered regions~\cite{abdelfattah2024coverage}. 
Our approach is complementary: by anchoring tests at root APIs and validating the resulting interaction trees via traces, we enable \emph{trace-based} coverage perspectives (e.g., diversity of downstream call paths per entry point), which we outline as opportunities in Section~\ref{sec:coverage-limitations}. In addition, availability-oriented validation can complement coverage views; the Service Availability Ratio~\cite{sar} quantifies microservice availability at the system level and offers an orthogonal quality dimension for test assessment.

\subsection{Positioning}
In summary, prior art either (i) targets single endpoints (model/search/fuzzing), (ii) enforces interface compatibility (CDC), (iii) replays multi-step workflows (E2E/grey-box), or (iv) diagnoses failures from traces post hoc. 
\emph{Root API Mode} synthesises these strengths while addressing their limitations: it confines execution to \emph{public entry points} (authentic client perspective), preserves ephemeral dependencies (no artificial step replays), and \emph{systematically} validates the full interaction chain through trace retrieval and analysis integrated into the testing/reporting pipeline.

%====================================================================
\section{Overview (Execution Modes: Overview and Selection)}\label{sec:overview}
This section describes the two execution modes implemented in the system and provides the rationale for selecting Root API Mode as the primary validation strategy. The definition of a step in Multi-Step Replay is clarified in Table \ref{tab:exec-modes}.

\begin{table}[t]
\centering
\caption{Execution Modes at a Glance}
\label{tab:exec-modes}
\renewcommand{\arraystretch}{1.15}\setlength{\tabcolsep}{4pt}
\begin{tabularx}{\columnwidth}{>{\bfseries}p{0.33\columnwidth} X}
\toprule
Mode & Summary \\
\midrule
Root API Mode (primary) & Executes publicly accessible entry points and validates the entire induced microservice interaction chain via post-execution analysis of distributed traces. Preserves authentic causality and ephemeral state without explicit workflow replay. \\
Multi-Step Replay (optional) & Explicitly executes each operation observed in a workflow trace as a separate client call, in dependency order. Intended for targeted reproduction and end-to-end checks when environments preserve cross-call prerequisites. \\
\bottomrule
\end{tabularx}
\end{table}

\textbf{Definition of a Step}:
A \emph{step} corresponds to a single inter-service operation observed in distributed tracing (typically one span for an HTTP call). Each step comprises: (i) operation signature (HTTP method and normalized path) and originating service; (ii) timing window and identifiers; (iii) extracted inputs (path, query, header, body) and observed outputs; (iv) an expected status derived from configuration with trace-informed fallback; and (v) parent/child links encoding causal ordering and dependencies. Steps compose into workflow scenarios whose partial order is induced by these links.

\subsection{Workflow Execution Challenges}\label{ssec:workflow-challenges}
While our system supports complete workflow execution for comprehensive end-to-end testing, practical deployment reveals fundamental challenges that make Root API Mode the preferred approach for microservice validation. Table~\ref{tab:workflow-challenges} summarizes these issues.

\begin{table*}[t]
\centering
\caption{Workflow Replication Challenges in Multi-Step Replay}
\label{tab:workflow-challenges}
\renewcommand{\arraystretch}{1.2}
\setlength{\tabcolsep}{6pt}
\begin{tabularx}{\textwidth}{>{\bfseries}p{0.22\textwidth} X p{0.26\textwidth}}
\toprule
Challenge & Description & Example / Impact \\
\midrule
Causal Chain Disruption & Executing Step~1 already triggers the full internal causal chain; downstream steps proceed naturally as part of the same request. & Replaying “Step~2” as a new client call is no longer on the original path; results lack causal continuity. \\
Duplicate Execution Artifacts & Re-invoking downstream steps as independent requests creates entirely new operations, not continuations. & Artificial execution patterns unlike real usage; skewed latency and coverage signals. \\
Ephemeral Data Loss & Traces omit ephemeral, in-memory values needed by downstream steps (temporary IDs, session state). & Replayed Step~2 lacks prerequisites set by Step~1, even though the original flow had them. \\
False Failure Oracles & Independent re-calls fail due to missing ephemeral dependencies, not true service faults. & Step~2 returns 4xx/5xx in replay despite end-to-end success in the original run; false negatives. \\
Misleading Causality & Failures in replayed Steps~2/3 appear as genuine service faults but are test artifacts. & Reports attribute errors to downstream services, obscuring true behavior and slowing diagnosis. \\
\bottomrule
\end{tabularx}
\end{table*}

\subsection{Coverage Analysis Limitations}\label{sec:coverage-limitations}

Traditional API testing coverage models face significant challenges when applied to our trace-driven single-API testing approach due to the paradigm shift from workflow-based to entry-point-focused test generation. This section details the current limitations and research opportunities in extending coverage analysis to this new testing methodology.

\paragraph{Single-Service Coverage Model Assumptions.}
RESTest's coverage framework follows established API testing coverage criteria \cite{martin2019restest}, including:
\begin{itemize}[leftmargin=*]
    \item \textit{Parameter Coverage:} Ensuring all input parameters are exercised
    \item \textit{Path Coverage:} Testing all API endpoints and URL templates  
    \item \textit{Operation Coverage:} Exercising all HTTP methods (GET, POST, etc.)
    \item \textit{Status Code Coverage:} Triggering different response codes
    \item \textit{Response Body Coverage:} Validating output schema elements
\end{itemize}

These criteria assume a unified OpenAPI specification and coherent parameter space, which breaks down in multi-service environments where each service has independent schemas and operation sets.

\paragraph{Trace-Driven Coverage Challenges.}
Our single-API testing approach with trace-based validation introduces novel coverage considerations:

% 1. Root API vs. Complete Path Coverage: While traditional coverage focuses on exercising all endpoints, our approach prioritizes comprehensive testing of entry-point APIs with validation of downstream service interactions through trace analysis.

1. Root API vs. Complete Path Coverage: Traditional coverage focuses on exercising all endpoints; this approach does not prioritize comprehensive testing of entry point APIs with validation of downstream service interactions through trace analysis.

2. Dynamic Service Chain Coverage: Coverage must account for variable service interaction patterns triggered by different parameter combinations, where the same API call may activate different microservice chains based on input values.

3. Trace-Based Validation Coverage: Traditional metrics cannot capture the coverage achieved through post-execution trace analysis, which validates service interactions that weren't directly tested but were triggered by root API calls.

4. AI-Enhanced Coverage Assessment: LLM-driven parameter generation creates coverage patterns that don't follow traditional combinatorial approaches, requiring new metrics to assess the semantic diversity of generated test cases.

\paragraph{Implementation Opportunities.}
Our trace-driven approach opens new directions for coverage analysis research:
\begin{itemize}[leftmargin=*]
    \item {Trace-Based Coverage Metrics:} Developing coverage models that measure the comprehensiveness of microservice interactions validated through distributed trace analysis
    \item {AI-Driven Coverage Assessment:} Creating metrics that evaluate the semantic diversity and fault-finding capability of LLM-generated test parameters
    \item {Dynamic Service Chain Coverage:} Measuring coverage across variable service interaction patterns triggered by different API parameter combinations
    \item {Failure Mode Coverage:} Assessing how well the test suite covers different types of microservice failure scenarios and error propagation patterns
\end{itemize}


\subsection{Why \textbf{Root API Mode} is Primary}\label{ssec:root-benefits}

The Root API Mode executes the publicly accessible API entry points and validates the entire induced microservice interaction chain through post-execution analysis of distributed traces. In this way, it preserves authentic causality and ephemeral state without explicit workflow replay. Additional benefits include:

\begin{itemize}[leftmargin=*]
  \item \emph{Authentic Testing:} Tests invoke only publicly accessible endpoints as external clients would, avoiding replay artifacts summarized in Table~\ref{tab:workflow-challenges}.
  \item \emph{Natural Execution Flow:} Ephemeral data dependencies and in-memory state transitions are preserved because downstream steps execute within the original causal chain.
  \item \emph{Comprehensive Validation:} Post-execution trace analysis validates the full service hierarchy triggered by a single API call, yielding broad coverage without artificial replays.
  \item \emph{AI-Enhanced Diagnosis:} Integration with LLMs provides intelligent root-cause analysis over the induced interaction chains.
\end{itemize}

Multi-Step Replay Mode offers complementary advantages to Root API Mode, but introduces redundant API calls that increase execution time and overhead. In contrast, Root API Mode optimizes performance by reducing unnecessary API interactions.


%====================================================================
\section{Methodology}\label{sec:method}
This section details the architectural and algorithmic design of the trace-driven extension for \textbf{RESTest}, which transforms distributed traces into targeted single-API test cases for comprehensive multi-service validation. The approach uses trace analysis to identify entry-point APIs, generates diverse test variants for these root endpoints, executes them individually, and conducts in-depth trace analysis to validate the complete microservice interaction chain. The design maintains RESTest's established \textit{generate → write → execute → report} pipeline, augmented with advanced trace-based orchestration and AI-driven failure analysis.
\begin{figure*}[t]
  \centering
  \includegraphics[width=\textwidth]{figs/Res_flow.png}
    \caption{Trace-driven Root API Mode testing pipeline. Inputs (\textit{blue}) consist of distributed traces and OpenAPI specifications. The processing flow identifies root APIs and maintains a registry (\textit{purple}), then generates tests through Smart Input Fetching and LLM-based synthesis (\textit{yellow}), writing executable cases (\textit{pink}). These are executed against the system (\textit{teal}), followed by trace retrieval and error analysis (\textit{red}). Outputs (\textit{green}) provide Allure reports with trace trees, validation summaries, and updates to the Root API Registry.}

  \label{fig:dataflow}
\end{figure*}
%====================================================================
\subsection{Tool Architecture}\label{ssec:arch}
The system comprises six specialized subsystems that collectively enable comprehensive microservice validation. While execution modes are defined in Section~\ref{sec:overview}, this subsection focuses on the architectural components supporting these modes. The subsystems include:

\begin{enumerate}[leftmargin=*]
  \item \textbf{Trace Workflow Extractor}  
        \hfill(\S\ref{ssec:extract})\\
        Parses OpenTelemetry~\cite{opentelemetry} or Jaeger~\cite{jaeger2017} span logs to extract complete workflow scenarios
        and identify root API endpoints that serve as natural system entry points (Algorithm~\ref{alg:extract}).
  \item \textbf{Test Generator}  
        \hfill(\S\ref{ssec:gen})\\
        Generates test cases in two modes: (i) \emph{Root API Mode}—single entry-point calls with
        trace-based post-execution validation; and (ii) \emph{Multi-Step Replay}—ordered inter-service
        steps with dependency-aware execution; both using LLM-driven parameter synthesis and
        smart input fetching for realistic test data (Algorithm~\ref{alg:generate}).
  \item \textbf{Multi-Service Test Writer}\\
        Emits JUnit/REST-assured test classes for both modes with
        integrated trace collection and comprehensive Allure reporting.
  \item \textbf{Trace-Aware Execution Engine}\\
        Executes tests (entry-point or multi-step) and captures trace identifiers. Integrates a short enhancement loop for failed cases. The loop fixes inputs and re-executes in rounds. It preserves reporting and analysis.
  \item \textbf{AI-Powered Trace Analysis Engine}\\
        Fetches execution traces, performs hierarchical failure analysis using
        \textit{TraceErrorAnalyzer}, and provides LLM-driven root cause diagnosis
        for the complete microservice interaction chains triggered by test executions.
  \item \textbf{Root API Registry System}  
        \hfill(\S\ref{ssec:registry})\\
        Maintains a persistent catalog of discovered root APIs with their complete
        microservice interaction patterns, enabling systematic test planning and
        long-term architectural analysis through intelligent deduplication.
\end{enumerate}




\noindent%
Figure~\ref{fig:dataflow} depicts the complete trace-driven testing pipeline, from initial trace extraction to AI-powered failure analysis. The diagram demonstrates the transformation of multi-service traces into focused single-API tests, ensuring comprehensive validation through post-execution trace analysis. Details regarding mode-specific generation and execution are provided in Section~\ref{ssec:gen-modes}.

%--------------------------------------------------------------------
\subsection{Trace-Driven Root API Extraction}
\label{ssec:extract}
\begin{algorithm}[tb]
\caption{TRACE\_WORKFLOW\_EXTRACTOR}\label{alg:extract}
\footnotesize
\begin{algorithmic}[1]
\Require\;{Path to trace file or directory}
\Ensure\;{List of \textsc{WorkflowScenario}s}
\State Parse JSON/JSONL files into span objects
\State $spansByTrace \gets$ Group spans by \texttt{traceId}
\ForAll{$(traceId, spans) \in spansByTrace$}
    \State $scenario \gets$ new \textsc{WorkflowScenario}
    \State $stepIndex \gets \{\}$ \Comment{spanId → WorkflowStep mapping}
    \ForAll{$span \in spans$}
        \State $step \gets$ \textproc{CreateWorkflowStep}$(span)$
        \State Extract HTTP method, path, headers, body from span tags
        \State Extract timing and status information
        \State $stepIndex[span.spanId] \gets step$
    \EndFor
    \State Link parent-child relationships using span references
    \State Identify root steps (spans with no parents)
    \State Add root steps to $scenario$
    \State Add $scenario$ to results
\EndFor
\State \textproc{MergeScenariosByDataDependency}(scenarios)
\State \Return scenarios
\end{algorithmic}
\end{algorithm}

\paragraph{Root API Identification and Accessibility Rationale.}
The \textproc{FindFirstBusinessAPI} function identifies the entry-point API call by filtering out authentication spans and selecting the first HTTP operation that represents actual business logic. This trace-driven approach addresses a fundamental challenge in microservice testing: \emph{API discoverability and accessibility}.

In microservice architectures, not all endpoints defined in OpenAPI specifications are directly accessible for testing. Many services implement \emph{internal} or \emph{private} APIs that are only accessible within the service mesh or cluster network, while others require complex authentication workflows or service-to-service communication protocols~\cite{waseem2020testing, zhou2018microservice}. A systematic review of microservice testing challenges identifies API visibility and service discoverability as primary obstacles to effective testing~\cite{microservice-testing-review-2023}.

By extracting root APIs from actual execution traces, we automatically identify the \emph{publicly accessible entry points} that real clients use to interact with the system. This approach ensures that our tests target APIs that are:
\begin{itemize}[leftmargin=*]
\item Actually accessible from external clients
\item Properly authenticated and authorized
\item Integrated with the system's API gateway or service mesh
\item Representative of real-world usage patterns
\end{itemize}

This methodology aligns with established best practices in microservice testing that emphasize focusing on externally accessible interfaces rather than attempting to test all internal service endpoints directly.

\paragraph{Trace Context Preservation.}
Although only the root API is used for test generation, the complete trace context is preserved to enable post-execution validation. This includes the full service interaction hierarchy, expected response patterns, and inter-service dependencies observed in the original execution.

%--------------------------------------------------------------------
\subsection{Test Generation}
\label{ssec:gen}

\subsubsection*{Generation Modes (summary)}\label{ssec:gen-modes}
Modes are introduced in Section~\ref{sec:overview}. Here we summarize their impact on generation only: Root API Mode produces single-call tests anchored at entry points with post-execution trace validation (see Section~\ref{ssec:exec}); Multi-Step Replay emits ordered step sequences with dependency-aware execution for environments that preserve cross-call prerequisites.

\begin{algorithm}[tb]
\caption{TEST\_GENERATOR}
\label{alg:generate}
\footnotesize
\begin{algorithmic}[1]
\Require\;{List of \textsc{WorkflowScenario}s, variant count $V$}
\Ensure\;{Collection of \textsc{MultiServiceTestCase}s}
\State Group scenarios by entry-point (HTTP method + normalized path)
\State Build shared parameter pools per group
\State $testCases \gets \{\}$
\ForAll{$scenario \in scenarios$}
    \For{$v \gets 1$ to $V$}
        \State $tc \gets$ new \textsc{MultiServiceTestCase}
        \State $context \gets \{\}$ \Comment{Runtime parameter context}
        \ForAll{$rootStep \in scenario.rootSteps$}
            \State \textproc{TraverseStep}$(rootStep, tc, context, v)$
        \EndFor
        \State Keep only the entry-point call in $tc$ \Comment{execution anchor}
        \State $testCases \gets testCases \cup \{tc\}$
    \EndFor
\EndFor
\State \Return $testCases$
\end{algorithmic}
\end{algorithm}

\paragraph{Hybrid parameter generation.}
\textproc{TraverseStep} applies a layered strategy that prioritizes realism and diversity:
\begin{itemize}[leftmargin=*]
  \item \textbf{System-resident values.} \emph{Smart Input Fetching} supplies existing IDs, codes, or tokens (used with configurable probability) to ensure valid baselines.
  % \item \textbf{Trace-derived hints.} When available, parameters are seeded from values observed in the original trace payloads.
  \item \textbf{LLM-based synthesis.} \emph{AiDrivenLLMGenerator} proposes semantically diverse candidates when fetching or trace seeding is insufficient.
\end{itemize}

\paragraph{Variants and efficiency.}
To balance coverage with cost, variants follow a consistent policy:
\begin{enumerate}[leftmargin=*]
  \item \textbf{High-fidelity baseline:} start from fetched/trace-seeded values to exercise realistic success paths.
  \item \textbf{Boundary exploration:} perturb constraints (ranges, enums, required/optional fields) to expose robustness issues.
  \item \textbf{Semantic diversity:} use LLM-generated values to probe alternative but plausible interpretations.
  \item \textbf{Targeted negatives:} inject purposeful violations to cover common failure modes and error handling.
\end{enumerate}

\paragraph{Context preservation.}
During generation, we maintain a runtime context for potential data dependencies and reuse across parameters; during execution, we anchor the call at the entry point and defer inter-service validation to post-execution trace analysis. This separation keeps execution authentic while enabling deep, trace-based oracles without replaying downstream steps for interservice communication.



\subsection{Test Case Types}\label{ssec:test-types}
The generator produces two case families with compact, configuration-driven oracles and AI support. \emph{Positive} cases exercise valid inputs and expect success; \emph{Negative} cases intentionally introduce invalid or semantically inconsistent inputs and expect errors.

\begin{table}[t]
\centering
\caption{Test Types and Validation Rules}
\label{tab:test-types}
\renewcommand{\arraystretch}{1.15}\setlength{\tabcolsep}{4pt}
\begin{tabularx}{\columnwidth}{>{\bfseries}p{0.24\columnwidth} X X}
\toprule
Type & Status-code rule & LLM response rule (2XX only) \\
\midrule
Positive & Pass only if \texttt{actual} $=$ \texttt{expectedStatus} (from configuration) & Fail if analysis reports a soft error; otherwise pass \\
Negative & Pass if \texttt{actual} $\neq$ \texttt{expectedStatus}; 2XX proceeds to soft-error check & Pass if analysis reports a soft error; otherwise fail \\
\bottomrule
\end{tabularx}
\end{table}
\paragraph{Negative case generation.}
\noindent This generation follows a four-stage pipeline: (i) \emph{Pool construction}—build per–root-API parameter pools from Smart Fetch values with LLM semantic expansion; (ii) \emph{Fault synthesis}—assemble typed invalids across eight categories (TYPE/REGEX/\textbf{SEMANTIC} mismatch, overflow, empty, null, special characters, boundary) using LLM; (iii) \emph{Variant selection}—controlled by \texttt{faulty.ratio} with round-robin (one invalid per test, cycling without repetition) or random (1--3 parameters, values may repeat); and (iv) \emph{Application \& oracles}—lock invalid values on selected parameters, apply normal Smart Fetch/LLM elsewhere, and validate via status code plus 2XX soft-error analysis.

AI is central to both generation and validation: LLMs synthesize semantically inconsistent inputs and detect \emph{soft errors} in successful responses, ensuring negative testing covers business‑logic semantics beyond surface constraints.

\paragraph{Selection strategies.}
Negative variants follow two configurable policies governed by concise properties: \texttt{faulty.ratio} controls the proportion of negatives; \texttt{faulty.round-robin} selects the strategy.
\begin{itemize}[leftmargin=*]
  \item \textbf{Round-robin}: one invalid parameter per test; cycles through all parameters and the eight fault types without repetition until pools are exhausted.
  \item \textbf{Random}: selects 1--3 parameters per test; fault types and values are chosen at random and may repeat.
\end{itemize}

\subsection{Root API Registry for Systematic Discovery and Management}\label{ssec:registry}

We introduce a concise \emph{Root API Registry} that persistently catalogues unique entry-point APIs (root APIs) together with their complete microservice interaction trees. The registry is populated automatically from distributed traces and used to plan tests, reason about dependency changes, and support regression targeting.

\paragraph{Design summary.}
The implementation consists of three lightweight components: (i) \textbf{RootApiRegistry}, a thread-safe manager that deduplicates root APIs and persists the registry; (ii) \textbf{RootApiEntry}, which represents a single root API and stores the set of distinct interaction trees observed for it; and (iii) \textbf{ApiTree}, a normalized tree that captures service call hierarchies (service, method, path) without request payloads. Trees are compared using structure-aware hashing to avoid storing duplicates while allowing multiple valid patterns per root API.

\paragraph{Applications.}
Distributed tracing has been shown to capture end-to-end causal relationships at scale, enabling dependency reconstruction and performance analysis (e.g., \textsc{Dapper}~\cite{sigelman2010dapper} and \textsc{Pivot Tracing}~\cite{mace2015pivot}). Empirical microservice benchmarks demonstrate that a single outward-facing API can exercise \emph{different internal call paths} under varying inputs and workloads, motivating the storage of multiple trees per root API (\textsc{DeathStarBench}~\cite{gan2020deathstarbench}). Systematic studies of microservice testing highlight persistent challenges in API discoverability, dependency awareness, and regression impact analysis, calling for trace-informed and automation-friendly artefacts~\cite{waseem2020testing}. Our registry operationalizes these insights by: (i) enumerating externally exercised entry points; (ii) recording the corresponding interaction trees; and (iii) persisting them for CI-aware test planning and change impact analysis.

\paragraph{Outputs and usage.}
The registry materializes two artefacts: (a) a de-duplicated list of root APIs (HTTP method + normalized path) discovered from traces; and (b) per-root \emph{sets of interaction trees} that document observed service-call hierarchies. The artefacts are stored in JSON for inspection and downstream analysis, and are updated incrementally across runs.

%--------------------------------------------------------------------
\subsection{Smart Input Fetching for Parameter Generation}\label{ssec:smartfetch}

The \emph{primary objective} of Smart Input Fetching is to harvest \emph{system-resident inputs} (e.g., existing identifiers, tokens, codes) so that requests use values that truly exist in the deployed environment. This prevents spurious client-side failures (e.g., \texttt{404} for non-existent IDs, \texttt{401} for missing tokens) and strengthens the oracle: we can observe whether an API succeeds or fails \emph{even with correct inputs}. Starting from valid seeds also enables purposeful negative tests by applying controlled mutations or withholding context values. The \emph{secondary objective} is efficiency and diversity: reduce expensive LLM calls while maintaining variety via semantic expansion.

\paragraph{Mechanism and integration.}
Configuration is supplied via the experiment’s properties (a feature toggle, a probability for using fetched values, the location of the registry, cache time-to-live, and the base service URL) and passed to the generator at startup. The generator reads these settings once, constructs a fetching configuration and fetcher, and consults an \emph{input fetch registry} that maps parameter patterns to provider endpoints with JSONPath-based extractors. Optional LLM-assisted discovery proposes new mappings from API metadata and service patterns. A cache bounds network overhead and stabilizes fetched values across a run; authentication hints can be provided for protected providers.


\paragraph{Testing intent and benefits.}
Using system-resident inputs makes tests representative and diagnostic: APIs are exercised under valid preconditions, revealing server-side faults that would be masked by invalid data. From these valid baselines, we construct deliberate negatives by perturbing selected fields, violating constraints, or withholding dependencies. The mechanism remains orthogonal to RESTest’s core pipeline: fetched values enter the same runtime context used for parameter injection and validation, while caching and timeouts keep cost predictable.


%--------------------------------------------------------------------
\subsection{Trace-Driven Test Execution and Post-Execution Analysis}
\label{ssec:exec}

The execution engine follows the standard lifecycle of test generation, writing, compilation, execution, and reporting. Tests are produced with integrated tracing and reporting hooks suitable for both Root API Mode and Multi-Step Replay. Algorithm~\ref{alg:execute} outlines the overall control flow, and Tables~\ref{tab:exec-overview}--\ref{tab:reporting-artefacts} summarize the execution pipeline, trace retrieval strategies, analysis modes, and reporting artefacts.

\begin{algorithm}[tb]
\caption{GENERATED\_TEST\_EXECUTION\_FLOW}\label{alg:execute}
\footnotesize
\begin{algorithmic}[1]
\Require\;{Generated JUnit test class with multiple test methods}
\Ensure\;{Allure reports with trace analysis}
\State Setup RestAssured with base URI and Allure filters
\ForAll{test method in generated class}
    \State Execute login step to obtain JWT token
    \State $stepResults \gets \{\}$, $failedSteps \gets 0$
    \ForAll{business step $s$ in test method}
        \State Start Allure step with metadata
        \State $startTime \gets$ current microseconds
        \State Build HTTP request with parameters and JWT
        \State $response \gets$ execute HTTP call
        \If{response status matches expected}
            \State $stepResults[s] \gets true$
            \State Attach response to Allure report
        \Else
            \State $stepResults[s] \gets false$, $failedSteps++$
            \State Attach detailed error information
        \EndIf
        \If{response status is 2XX \textbf{and} soft-error validation enabled}
            \State $vr \gets$ \textproc{LLMResponseValidator.validate}(service, endpoint, status, body)
            \State Attach \emph{INTELLIGENT ANALYSIS (Based on Response)} with RCA
            \State Update $stepResults[s]$ according to test type (positive expects SUCCESS; negative expects FAILED)
            \If{verdict changes from pass to fail}
                \State $failedSteps++$
            \EndIf
        \EndIf
        \State $traces \gets$ \textproc{FetchJaegerTraces}$(service, method, path, startTime)$
        \If{traces found}
            \State $analysis \gets$ \textproc{TraceErrorAnalyzer.analyze}$(traces)$
            \State Generate hierarchical trace visualization
            \State Attach trace tree and analysis to Allure
            \If{$analysis.hasErrors()$}
                \State $llmDiagnosis \gets$ \textproc{LLMRootCauseAnalysis}$(analysis)$
                \State Attach AI-generated failure diagnosis
            \EndIf
        \EndIf
        \State Stop Allure step
    \EndFor
    \State Generate comprehensive test summary
    \If{$failedSteps > 0$ OR login failed}
        \State Mark test as FAILED with detailed breakdown
    \Else
        \State Mark test as PASSED
    \EndIf
\EndFor
\end{algorithmic}
\end{algorithm}

\paragraph{Soft-error detection (2XX).}
Soft errors denote business-logic failures surfaced in the response body despite successful HTTP statuses. When enabled, the validator inspects explicit failure flags (e.g., \texttt{status:0}, \texttt{success:false}), error fields (\texttt{error}, \texttt{msg}, \texttt{message}), null or empty \texttt{data} when payloads are expected, and constraint/validation messages. The validator is implemented with large language models that reason over JSON responses to detect \emph{semantic} failures beyond status codes and schemas: it interprets business indicators and messages (including varied phrasing/locales), checks payload consistency (e.g., expected \texttt{data} present and non-empty), and recognises domain-level invariants and constraint violations. Prompts emphasise precision to avoid false positives; evaluation is limited to 2XX responses and produces a concise RCA attached to the report. Positive tests fail when a soft error is detected; negative tests pass when a soft error is detected in a 2XX response.

\paragraph{Execution Overview.}
To facilitate interpretation, Table~\ref{tab:exec-overview} presents the principal execution stages together with their inputs and outputs. Trace acquisition policies are consolidated in Table~\ref{tab:trace-retrieval}. Diagnostic capabilities are organized in Table~\ref{tab:analysis-modes}. Finally, Table~\ref{tab:reporting-artefacts} enumerates the reporting artefacts made available to developers and analysts.

\begin{table*}[t]
\centering
\caption{Execution Pipeline Overview}
\label{tab:exec-overview}
\renewcommand{\arraystretch}{1.15}\setlength{\tabcolsep}{6pt}
\begin{tabularx}{\textwidth}{>{\bfseries}p{0.17\textwidth} X p{0.26\textwidth}}
\toprule
Phase & Key Actions & Outputs \\
\midrule
Setup & Configure RestAssured; register Allure filters; initialize runtime context. & Base URI configured; request/response interceptors active. \\
Authentication & Obtain JWT/token as Step~0 when required; cache for subsequent calls. & Valid auth headers for requests. \\
Step Execution & Build requests with synthesized parameters; execute HTTP calls; record timing. & Responses with bodies, headers, status; timing metrics. \\
Validation & Status-code check against configuration; \emph{plus} 2XX soft-error detection via LLM when enabled. & Per-step verdicts updated by soft-error analysis; error details; RCA attachments. \\
Trace Capture & Query Jaeger/OTel using narrow time window and multi-strategy search; identify API spans. & Retrieved trace with span hierarchy. \\
Analysis & Run TraceErrorAnalyzer; optionally invoke LLM for diagnostics. & Root-cause findings; suggested fixes. \\
Reporting & Emit Allure steps, trace trees, and summaries. & Rich report with artefacts for debugging. \\
\bottomrule
\end{tabularx}
\end{table*}

\begin{table*}[t]
\centering
\caption{Trace Retrieval Strategies}
\label{tab:trace-retrieval}
\renewcommand{\arraystretch}{1.15}\setlength{\tabcolsep}{6pt}
\begin{tabularx}{\textwidth}{>{\bfseries}p{0.24\textwidth} X p{0.24\textwidth}}
\toprule
Strategy & Description & When Used \\
\midrule
Precise Time Window & Query from 10s before to 30s after API start to localize trace ingest latency. & Default for all calls; narrows candidate traces. \\
Multi-Strategy Queries & Combine operation-match, service-based, and time-range searches to improve recall. & When operation names vary or services rename endpoints. \\
Retry with Backoff & Re-issue queries with exponential backoff to tolerate ingestion delays. & When initial queries miss late-arriving spans. \\
API Call Detection & Filter spans to keep real API calls vs. internal framework operations. & Always; to avoid noise in analysis. \\
\bottomrule
\end{tabularx}
\end{table*}

\paragraph{Trace-centred diagnosis.}
After trace capture, analysis operates on the span graph to reveal technical root causes and propagation paths. The process identifies error spans, reconstructs parent--child relationships, isolates minimal root-cause sets, and extracts evidence (exceptions, stack traces, status codes, timings) for reporting and remediation.

\begin{table*}[t]
\centering
\caption{Trace Error Analysis Workflow}
\label{tab:trace-workflow}
\renewcommand{\arraystretch}{1.15}\setlength{\tabcolsep}{6pt}
\begin{tabularx}{\textwidth}{>{\bfseries}p{0.28\textwidth} X p{0.26\textwidth}}
\toprule
Step & Key Action & Outputs \\
\midrule
Error span identification & Scan spans flagged as errors vs. expected business failures & Candidate error set \\
Hierarchy mapping & Build parent--child call graph and propagation paths & Causal chain per error \\
Root-cause isolation & Find errors with no failed children (true sources) & Minimal root-cause set \\
Evidence extraction & Parse logs/exceptions, status, timings & Failure details for RCA \\
\bottomrule
\end{tabularx}
\end{table*}

\begin{table*}[t]
\centering
\caption{Failure Analysis Modes}
\label{tab:analysis-modes}
\renewcommand{\arraystretch}{1.15}\setlength{\tabcolsep}{6pt}
\begin{tabularx}{\textwidth}{>{\bfseries}p{0.24\textwidth} X p{0.24\textwidth}}
\toprule
Mode & Capabilities & Outputs \\
\midrule
TraceErrorAnalyzer (rule-based) & Identify failed spans; map parent/child propagation; isolate root-cause spans; extract stack traces and metadata. & Root-cause spans; failure locations; propagation summary. \\
LLM Response Validation (2XX soft errors) & Detect business-logic failures in successful responses; produce RCA based on body content. & Soft-error verdict; RCA attachment; updated test outcome. \\
LLM-Enhanced Diagnosis (optional) & Generate technical explanations, remediation suggestions, and dependency impact given trace context. & Human-readable diagnosis; actionable fixes; impact assessment. \\
\bottomrule
\end{tabularx}
\end{table*}

\paragraph{LLM-assisted diagnosis and fallback.}
Diagnosis leverages models to translate trace context into clear technical explanations and remediation guidance, including dependency impact. When models are unavailable, a rule-based analyser classifies exceptions, recognises common failure patterns, pinpoints failure locations from stack traces, and outlines likely service impacts.

\begin{table*}[t]
\centering
\caption{Reporting Artefacts (Allure)}
\label{tab:reporting-artefacts}
\renewcommand{\arraystretch}{1.15}\setlength{\tabcolsep}{6pt}
\begin{tabularx}{\textwidth}{>{\bfseries}p{0.24\textwidth} X p{0.24\textwidth}}
\toprule
Artefact & Contents & Purpose \\
\midrule
Allure Steps & Per-step metadata (service, method, path, expected status), request/response attachments, timing. & Reproducibility; quick failure localization. \\
Trace Visualization & Hierarchical span tree with timing/status and service edges. & Understand end-to-end behavior for each call. \\
INTELLIGENT ANALYSIS (Based on Response) & LLM soft-error detection for 2XX bodies with concise RCA. & Validate business-logic correctness when status is 2XX. \\
AI Diagnostics (Trace) & LLM-generated root-cause analysis from traces (when enabled). & Accelerate debugging and remediation. \\
Service Health Metrics & Aggregated timings and error rates across invoked services. & Spot hotspots and regressions. \\
\bottomrule
\end{tabularx}
\end{table*}

\paragraph{Post-Execution Input Enhancement}
After initial runs, the execution pipeline performs a focused input enhancement stage. The stage collects failed cases with a runtime listener. The stage filters server-side errors when configured. The system prepares compact contexts that include endpoint, method, status, parameters, and failure message. The system then requests targeted input fixes from the model.

The pipeline regenerates test files with minimal edits. The pipeline recompiles and re-executes in short rounds. Reports mark enhanced cases and attach both the original failure and revised inputs. Outputs persist per test and per round. The loop keeps changes explainable and avoids masking backend faults.

Prompts follow a strict structure for clarity and control. Each prompt summarizes endpoint and method, encodes status and failure message, and enumerates parameters with types and locations. The request asks for improved values that satisfy constraints and preserve semantics. The model returns enhanced values with concise reasoning. The system applies only changes that address the reported issue and preserves other fields.

\noindent The writer also creates timestamped test packages to isolate runs; reports can be aggregated for trends while preserving per-run debugging. Coverage analysis is adapted to single-API testing patterns (Section~\ref{sec:coverage-limitations}).





%====================================================================
\section{Experiments}\label{sec:experiments}

\subsection{System Under Test}
We evaluate our Root API Mode approach on \emph{TrainTicket}~\cite{zhou2018trainticket}, a representative microservice architecture comprising multiple interconnected services. Our system extracts root API endpoints from distributed traces spanning ts-admin-basic-info-service, ts-price-service, ts-admin-order-service, ts-station-service, ts-route-service, and ts-admin-route-service. The Root API Mode approach focuses testing on externally accessible entry points while validating comprehensive microservice interactions through post-execution trace analysis, avoiding the workflow replication challenges described in Section~\ref{ssec:workflow-challenges}.

\subsection{Root API Extraction and Test Generation}
We collect five representative traces from TrainTicket's web frontend, each capturing a distinct user operation and its complete microservice interaction chain. Operating in Root API Mode, we extract the externally accessible entry points (excluding internal authentication flows) as primary test targets. The Root API Registry automatically catalogues these discoveries, identifying 8 unique root APIs with 15 distinct microservice interaction patterns. For each root API endpoint, we generate ten diverse test variants using our hybrid parameter generation strategy:

\textbf{Smart Input Fetching:} Primary variants use the \emph{SmartInputFetcher} to harvest valid system values, ensuring realistic baseline tests that reflect authentic client interactions.

\textbf{LLM-Driven Exploration:} Additional variants employ the \emph{AiDrivenLLMGenerator} with semantic expansion to explore parameter boundaries and edge cases while maintaining system compatibility.

\textbf{Root API Focus:} The 50\% probability split between smart-fetched and LLM-generated values ensures that tests exercise root APIs with both realistic and exploratory data, avoiding the ephemeral dependency issues inherent in workflow replication.

\subsection{Root API Mode Execution and Trace Analysis Protocol}
Our Root API Mode execution protocol emphasizes natural client interaction patterns with comprehensive validation:

\textbf{Entry Point Execution:} Each test executes only the root API endpoint as external clients would, allowing natural microservice causal chains to execute organically within the system.

\textbf{Trace Capture:} We capture trace identifiers during API execution and fetch complete distributed traces using multiple query strategies, typically waiting 2-5 seconds for comprehensive trace collection.

\textbf{Comprehensive Analysis:} The \textsc{TraceErrorAnalyzer} examines the complete service interaction tree triggered by root API calls, while LLM integration provides intelligent diagnosis for complex distributed failures.

\textbf{Authentic Reporting:} Enhanced Allure reports visualize the natural microservice interaction hierarchies triggered by root API executions, providing AI-generated failure analysis and debugging recommendations based on authentic system behavior.

\textbf{Post-Execution Enhancement:} For failed cases, we fix inputs with model suggestions and re-run in short rounds. Reports mark enhanced cases and attach both failure and new inputs. The step respects skip rules for 5XX errors.

\subsection{Results and Trace-Based Evidence}
Figure~\ref{fig:case_exec} demonstrates our approach's capability to detect backend failures through single API execution. The test targets POST /api\slash v1\slash adminbasicservice\slash adminbasic\slash prices on ts-admin-basic-info-service with LLM-generated parameters. While the API call itself returns HTTP 500, our trace analysis reveals the complete failure propagation through the microservice chain. The enhanced Allure report shows not only the immediate API failure but also the subsequent trace retrieval and analysis, including AI-generated root cause diagnosis that identifies the specific service component responsible for the failure.

\begin{figure}[!t]
\centering
\includegraphics[width=\linewidth]{case_exec.png}
\caption{End to end execution of one generated case. Login succeeds. The call to /api\slash v1\slash adminbasicservice\slash adminbasic\slash prices returns 500. The follow up call to /api\slash v1\slash priceservice\slash prices does not meet the expected 201. Allure stores request bodies, requests, responses, and error details.}
\label{fig:case_exec}
\end{figure}

Figure~\ref{fig:post_prices_matrix} shows the comprehensive test matrix for our root API approach targeting POST /api\slash v1\slash adminbasicservice\slash adminbasic\slash prices. The matrix displays ten test variants generated through our hybrid parameter strategy. Two tests achieve successful execution (HTTP 200/201), while several variants trigger \emph{Status Code Mismatch} failures. Each test result includes trace analysis showing whether failures originate from the target service or propagate from downstream dependencies. The matrix demonstrates how parameter variations expose different failure modes in the microservice architecture.

\begin{figure}[!t]
\centering
\includegraphics[width=\linewidth]{post_prices_matrix.png}
\caption{Per call results for POST /api\slash v1\slash adminbasicservice\slash adminbasic\slash prices. Two tests succeed. Several tests fail with \emph{Status Code Mismatch}. The runner holds a valid token in all runs.}
\label{fig:post_prices_matrix}
\end{figure}

Figure~\ref{fig:get_adminorder_matrix} shows another API. The call is GET /api\slash v1\slash adminorderservice\slash adminorder. All shown runs return 200. One run takes about 1454 ms. Most runs complete near 20 to 64 ms. The grid gives a quick scan for both correctness and latency spread under new inputs.

\begin{figure}[!t]
\centering
\includegraphics[width=\linewidth]{get_adminorder_matrix.png}
\caption{Per call results for GET /api\slash v1\slash adminorderservice\slash adminorder. All runs return 200. One run shows a higher latency spike.}
\label{fig:get_adminorder_matrix}
\end{figure}

Figure~\ref{fig:adminroute_exec} demonstrates our trace-driven analysis of a complex microservice failure. The test executes a single API call to POST /api\slash v1\slash adminrouteservice\slash adminroute on ts-admin-route-service, which returns HTTP 400. However, our post-execution trace analysis reveals the complete service interaction chain that would have been triggered by a successful call, including subsequent calls to ts-station-service and ts-route-service. The enhanced Allure report provides AI-generated analysis identifying the exact validation failure in the admin route service, along with debugging recommendations. This approach enables developers to understand both the immediate failure and its potential impact on the broader microservice ecosystem.

\begin{figure}[!t]
\centering
\includegraphics[width=\linewidth]{adminroute_exec.png}
\caption{Three step workflow that includes ts-admin-route-service, ts-station-service, and ts-route-service. The first business step returns 400. The following steps return 200. The scenario fails due to step one.}
\label{fig:adminroute_exec}
\end{figure}

Figure~\ref{fig:smartfetch_body} demonstrates our Smart Input Fetching capability in action. The figure shows a request body populated with actual station names (haidian, guangzhou, dalian) harvested from the running TrainTicket system. This realistic data generation ensures that tests use valid system references rather than synthetic values, leading to more authentic failure scenarios. When this smart-fetched data is used in API calls, it enables the complete microservice chain to execute successfully, providing comprehensive validation of service interactions and performance characteristics.

\begin{figure}[!t]
\centering
\includegraphics[width=\linewidth]{smartfetch_body.png}
\caption{Smart Fetch example for POST /api\slash v1\slash stationservice\slash stations\slash idlist. The generator fills the body with existing station names. The service returns 200.}
\label{fig:smartfetch_body}
\end{figure}

\subsection{Analysis}
Our Root API Mode approach demonstrates superior fault detection and diagnosis capabilities compared to traditional workflow replication. The focus on externally accessible entry points eliminates the false failure oracles and ephemeral dependency issues that plague workflow-based testing.

\textbf{Authentic Failure Detection:} Root API calls trigger natural microservice failure cascades, as evidenced by the price service HTTP 500 errors and admin route validation failures. These represent genuine system issues rather than testing artifacts created by workflow replication.

\textbf{Natural Execution Patterns:} By invoking only the APIs that external clients would naturally call, our approach preserves in-memory data dependencies and ephemeral state transitions, ensuring that trace analysis reflects authentic system behavior.

\textbf{Comprehensive Validation:} Post-execution trace analysis captures the complete service interaction hierarchies triggered by root API executions, providing deeper validation coverage than artificial step-by-step workflow replication.

\textbf{Superior Debugging:} Enhanced Allure reports visualize authentic microservice interaction patterns with AI-generated root cause analysis. The elimination of false causality chains significantly improves debugging accuracy and reduces time-to-resolution for genuine system failures.

\subsection{Takeaways}
Our Root API Mode testing approach achieves comprehensive microservice validation while avoiding the fundamental challenges of workflow replication:

\textbf{Authentic System Testing:} Root API Mode tests exercise microservice systems exactly as external clients would, preserving natural execution flows and ephemeral data dependencies that workflow replication cannot capture.

\textbf{Elimination of False Oracles:} By avoiding artificial step-by-step replay, the approach eliminates misleading test failures caused by missing ephemeral dependencies, providing reliable validation signals.

\textbf{Enhanced Fault Detection:} The method successfully exposes genuine backend failures through natural API invocation patterns, while post-execution trace analysis provides comprehensive validation of the complete microservice interaction chains.

\textbf{AI-Enhanced Diagnosis:} Integration of LLM-powered failure analysis with authentic trace data transforms complex distributed failures into actionable debugging information based on real system behavior.

\textbf{Practical Deployment Value:} Root API Mode provides superior practical value for microservice validation by focusing on the natural system boundaries and interaction patterns that matter for production deployments.

\textbf{Systematic Knowledge Accumulation:} The Root API Registry transforms discovery-based testing into a systematic process where each test execution contributes to a growing knowledge base of microservice interaction patterns, enabling long-term architectural analysis and comprehensive test planning.



\bibliographystyle{IEEEtran}
% TODO: add references.bib once experiments are finalised.
\bibliography{references}
\end{document}
